% Generated draft: requires mathematical and source review.
\documentclass[11pt]{article}
\usepackage[margin=25mm]{geometry}
\usepackage{fontspec}
\setmainfont{Latin Modern Roman}
\usepackage{amsmath,amssymb,mathtools,unicode-math}
\setmathfont{Latin Modern Math}
\usepackage{enumitem,needspace,xurl,hyperref}
\setlength{\parindent}{0pt}
\setlength{\parskip}{5pt}
\setlength{\emergencystretch}{3em}
\newcommand{\N}{\mathbb{N}}
\newcommand{\Z}{\mathbb{Z}}
\newcommand{\Q}{\mathbb{Q}}
\newcommand{\R}{\mathbb{R}}
\newcommand{\C}{\mathbb{C}}
\newcommand{\F}{\mathbb{F}}
\newcommand{\A}{\mathbb{A}}
\newcommand{\PP}{\mathbb{P}}
\newcommand{\E}{\mathbb{E}}
\newcommand{\eps}{\varepsilon}
\newcommand{\dd}{\,\mathrm{d}}
\begin{document}
\section*{UnsolvedMath ID: 20003104}
\textbf{Source problem code:} catalogue entry
\subsection{Definitions and mathematical statement}
A. Papadopoulos asks for a good symmetrization of the Thurston metric; the source notes that Sorvali's length-spectrum metric is one. In the usual finite-type Teichm\"uller space with nonzero dimension, fix the normalization
\[
d_{\mathrm{Th}}(X,Y)=\log\sup_{\alpha}\frac{\ell_Y(\alpha)}{\ell_X(\alpha)},\qquad
d_{\mathrm{ls}}(X,Y)=\max\{d_{\mathrm{Th}}(X,Y),d_{\mathrm{Th}}(Y,X)\}.
\]
Thus $d_{\mathrm{ls}}=\sup_\alpha|\log\ell_Y(\alpha)-\log\ell_X(\alpha)|$, with no factor $1/2$ in this convention. The curve set consists of essential simple closed classes. For a basepoint $X_0$, the normalized log-length map
\[
\Phi(X)=\left(\log\frac{\ell_X(\alpha)}{\ell_{X_0}(\alpha)}\right)_\alpha\in\ell^\infty(\mathcal C(S))
\]
realizes $d_{\mathrm{ls}}$ as the ambient chord metric $\|\Phi(X)-\Phi(Y)\|_\infty$.

The aid contrasts this with infinitesimal symmetrization. Let
\[
F_X(v)=\sup_\alpha\dd(\log\ell_\alpha)_X(v),\qquad
F_X^\vee(v)=\max\{F_X(v),F_X(-v)\}
=\sup_\alpha|\dd(\log\ell_\alpha)_X(v)|.
\]
Its intrinsic path metric is
\[
\delta_\vee(X,Y)=\inf_{\gamma:X\to Y}\int F_{\gamma(t)}^\vee(\dot\gamma(t))\,\dd t,
\]
using piecewise smooth, or the corresponding absolutely continuous, paths. The aid asserts that this is the intrinsic length metric of the log-length image, is reversible, mapping-class invariant, complete, topology-compatible and geodesic, and satisfies
\[
d_{\mathrm{ls}}(X,Y)\leq\delta_\vee(X,Y).
\]
The surface class and analytic hypotheses for these metric properties need the full source argument; they are not extended here to arbitrary infinite-type spaces.

For a forward-dominant pair, meaning $D=d_{\mathrm{Th}}(X,Y)\geq d_{\mathrm{Th}}(Y,X)$, the aid's calibrated criterion is
\[
\delta_\vee(X,Y)=d_{\mathrm{ls}}(X,Y)=D
\]
if and only if there exists a forward Thurston geodesic $\gamma$ from $X$ to $Y$ such that
\[
F_{\gamma(t)}(\dot\gamma(t))\geq F_{\gamma(t)}(-\dot\gamma(t))
\quad\text{for almost every }t.
\]
The path is a minimizing $\delta_\vee$ path as well; arclength reparametrization is allowed. If every forward Thurston geodesic between the pair violates this inequality on a positive-measure set, then $\delta_\vee(X,Y)>d_{\mathrm{ls}}(X,Y)$. The reverse-dominant case is handled by reversing the ordered pair. This is a pairwise conditional test, not an exhibited strict pair or a proof of equality for all pairs. ``Good'' is not given a unique optimization criterion in the original problem, so these constructions and comparisons do not select a universally preferred symmetrization.
\subsection{Short English statement}
Compare the usual endpoint length-spectrum symmetrization with the intrinsic metric obtained by symmetrizing Thurston's infinitesimal norm. The supplied aid gives their inequality and an exact geodesic sign condition for equality, rather than claiming they always coincide.
\subsection{Sources}
\begin{itemize}
\item Local UnsolvedMath registry, permanent ID $20003104$, source code \texttt{catalogue entry}; canonical \texttt{aim-topology-notes.json}, \texttt{notes[191]}.
\item AIM \emph{Problems on Thurston Metric}, Section 3, Problem 3.1, attributed to A. Papadopoulos, workshop \emph{Lipschitz metric on Teichmueller space}: \url{https://aimath.org/pastworkshops/lipschitzteichproblems.pdf}. The source names Sorvali's length-spectrum metric; original remarks and literature field are empty. Trailing dots are extraction punctuation.
\item Research attempt 1: \texttt{partial\_result}, status \texttt{partially\_solved}, classification \texttt{PARTIAL-PROGRESS}, rigor \texttt{proved}, novelty confidence low. No authoritative complete references for the intrinsic-metric assertions are supplied. This machine-generated aid is not peer reviewed and requires independent verification before citation.
\end{itemize}
\end{document}
